
\documentclass[12pt,a4paper]{article}

% =========================
% PACKAGES
% =========================
\usepackage{amsmath}
\usepackage{amssymb}
\usepackage{amsthm}
\usepackage{geometry}
\usepackage{enumitem}
\usepackage{mathtools}

\geometry{margin=1in}

% =========================
% DOCUMENT INFORMATION
% =========================
\title{Mathematical LaTeX Practice}
\author{Shibam Mallick}
\date{\today}

% =========================
% THEOREM ENVIRONMENTS
% =========================
\newtheorem{theorem}{Theorem}
\newtheorem{definition}{Definition}
\newtheorem{example}{Example}

% =========================
% DOCUMENT
% =========================
\begin{document}

\maketitle

\tableofcontents

\newpage

% =========================
% 1. BASIC TEXT
% =========================
\section{Introduction}

This is a practice document for learning \LaTeX.

We can write mathematics both inside a sentence, such as
\(x^2 + y^2 = z^2\), and as a separate displayed equation:

\[
x^2 + y^2 = z^2.
\]

% =========================
% 2. BASIC MATHEMATICS
% =========================
\section{Basic Mathematics}

\subsection{Fractions}

A fraction can be written as

\[
\frac{a}{b}.
\]

A more complicated fraction is

\[
\frac{x^2 + 2x + 1}{x^2 - 1}.
\]

\subsection{Powers and Roots}

\[
x^2,\qquad
x^n,\qquad
\sqrt{x},\qquad
\sqrt[n]{x}.
\]

For example,

\[
\sqrt{x^2 + y^2}.
\]

% =========================
% 3. EQUATIONS
% =========================
\section{Equations}

The quadratic equation is

\[
ax^2 + bx + c = 0.
\]

Its solutions are

\[
x = \frac{-b \pm \sqrt{b^2 - 4ac}}{2a}.
\]

Multiple equations can be aligned using the \texttt{align} environment:

\begin{align}
2x + 3 &= 7 \\
2x &= 4 \\
x &= 2.
\end{align}

% =========================
% 4. SETS
% =========================
\section{Set Theory}

Let

\[
A = \{1,2,3,4,5\}
\]

and

\[
B = \{3,4,5,6,7\}.
\]

Then

\[
A \cup B = \{1,2,3,4,5,6,7\},
\]

\[
A \cap B = \{3,4,5\}.
\]

The empty set is

\[
\varnothing.
\]

% =========================
% 5. SUMMATION AND PRODUCTS
% =========================
\section{Summation and Products}

The sum of the first \(n\) positive integers is

\[
\sum_{i=1}^{n} i
=
\frac{n(n+1)}{2}.
\]

A product can be written as

\[
\prod_{i=1}^{n} i = n!.
\]

% =========================
% 6. LIMITS
% =========================
\section{Limits}

Consider the limit

\[
\lim_{x\to 0}\frac{\sin x}{x}=1.
\]

Another example is

\[
\lim_{x\to\infty}\frac{1}{x}=0.
\]

% =========================
% 7. DIFFERENTIATION
% =========================
\section{Differentiation}

Let

\[
f(x) = x^3 + 2x^2 - 5x + 1.
\]

Then

\[
f'(x)
=
3x^2 + 4x - 5.
\]

The second derivative is

\[
f''(x)
=
6x + 4.
\]

% =========================
% 8. INTEGRATION
% =========================
\section{Integration}

An indefinite integral is

\[
\int x^2\,dx
=
\frac{x^3}{3}+C.
\]

A definite integral is

\[
\int_0^1 x^2\,dx
=
\frac{1}{3}.
\]

% =========================
% 9. MATRICES
% =========================
\section{Matrices}

A \(2\times2\) matrix can be written as

\[
A =
\begin{pmatrix}
a & b \\
c & d
\end{pmatrix}.
\]

Its determinant is

\[
\det(A)
=
\begin{vmatrix}
a & b \\
c & d
\end{vmatrix}
=
ad-bc.
\]

% =========================
% 10. VECTORS
% =========================
\section{Vectors}

A vector can be written as

\[
\mathbf{v}
=
\begin{pmatrix}
x\\
y\\
z
\end{pmatrix}.
\]

The magnitude of the vector is

\[
\|\mathbf{v}\|
=
\sqrt{x^2+y^2+z^2}.
\]

% =========================
% 11. PROBABILITY
% =========================
\section{Probability}

The probability of an event \(A\) is

\[
P(A)
=
\frac{\text{Number of favourable outcomes}}
{\text{Total number of outcomes}}.
\]

For two events,

\[
P(A\cup B)
=
P(A)+P(B)-P(A\cap B).
\]

If \(A\) and \(B\) are independent, then

\[
P(A\cap B)
=
P(A)P(B).
\]

% =========================
% 12. THEOREM
% =========================
\section{Theorem}

\begin{theorem}
If \(f'(x)>0\) on an interval \(I\), then \(f\) is strictly increasing on \(I\).
\end{theorem}

\begin{proof}
Let \(x_1,x_2\in I\) with \(x_1<x_2\).
By the Mean Value Theorem, there exists
\(c\in(x_1,x_2)\) such that

\[
f(x_2)-f(x_1)
=
f'(c)(x_2-x_1).
\]

Since

\[
f'(c)>0
\]

and

\[
x_2-x_1>0,
\]

we have

\[
f(x_2)-f(x_1)>0.
\]

Therefore,

\[
f(x_2)>f(x_1).
\]

Hence \(f\) is strictly increasing on \(I\).
\end{proof}

% =========================
% 13. LISTS
% =========================
\section{Important Topics}

\begin{enumerate}
    \item Algebra
    \item Calculus
    \item Linear Algebra
    \item Probability
    \item Statistics
\end{enumerate}

Some important mathematical symbols are:

\begin{itemize}
    \item \(\in\) -- belongs to
    \item \(\subseteq\) -- subset
    \item \(\forall\) -- for all
    \item \(\exists\) -- there exists
    \item \(\Rightarrow\) -- implies
    \item \(\Leftrightarrow\) -- if and only if
\end{itemize}

% =========================
% 14. EXAMPLE
% =========================
\section{Example}

\begin{example}
Find the derivative of

\[
f(x)=x^4-3x^2+2.
\]

Using the power rule,

\[
f'(x)
=
4x^3-6x.
\]

Therefore,

\[
\boxed{f'(x)=4x^3-6x}.
\]
\end{example}

% =========================
% END
% =========================

\end{document}
```
